\documentclass[runningheads]{llncs}
\usepackage[T1]{fontenc}
\usepackage{lmodern,microtype,amsmath,amssymb,booktabs,graphicx,tabularx,array,longtable}
\usepackage[hidelinks]{hyperref}
\graphicspath{{../artifacts/figures/}}
\input{generated/macros}
\input{generated/semantic/macros}
\newcommand{\method}{M1}
\newcommand{\suppTable}[3]{\begin{table}[!ht]\caption{#2}\label{#3}\centering\scriptsize
\setlength{\tabcolsep}{3pt}\input{#1}\end{table}}
\begin{document}
\title{Supplement: Semantic Structure and Revision in Temporal Evidence Graphs for Generative Monitoring}
\titlerunning{Temporal Evidence Graphs: Supplement}
\author{Anonymous Authors}\authorrunning{Anonymous Authors}
\institute{Anonymous Institution}\maketitle
\setcounter{section}{0}
\renewcommand{\thesection}{S\arabic{section}}
\renewcommand{\thetable}{S\arabic{table}}
\renewcommand{\thefigure}{S\arabic{figure}}

\section{Construction and Projection}
\input{generated/semantic/full_nodes}
\input{generated/semantic/full_relations}

The node table covers all seven Record kinds plus adapter-created assessments and
unresolved reference placeholders. The latter are possible through source-endpoint
MERGE operations but have zero observed count in the exported original scopes.
Domains describe adapter behavior, not a narrower schema that the database does
not enforce. In particular, source IDs produce broad Record-to-Record relations;
admission checks supply the numerical and task restrictions. Identity constraints
enforce scope/ID uniqueness. Acyclicity and immutable-record checks are implemented
in Python; neither OWL entailment nor SHACL validation is claimed.

The conceptual PROV-DM mapping treats observation and feature versions as entities,
feature extraction and generation as activities recorded in metadata, and
source/version links as derivation and revision information. The actual database
does not materialize a complete PROV activity vocabulary. A SupportAssessment is
an outcome keyed by scope, record ID, knowledge time and trigger ID. The SQLite
control stores equivalent assessment rows without Neo4j's extra assessment nodes.
CONTRADICTS has no edge timestamp, so its associated assessment supplies time.
Raw rejected proposals are archived separately from persisted ClaimVersions.

An audit finds \SemAssessmentTies{} record/time groups with multiple assessments,
including \SemAssessmentConflicts{} with different states across the original
exported scopes. Sequential triggers can share an ingestion timestamp without
an explicit within-time sequence field. The semantic projection preserves every
tied outcome and marks conflicting summaries ambiguous, rather than treating
database element order as chronology. Independent proposition scores and display
membership use the event log and saved explanation versions, so this audit
changes neither the original outcomes nor the new grounding/recall results.

The stored graph contains immutable versions, explanation containers, subject and
sensor identities, review events and assessments. The active evidence-and-claim
graph selects latest visible logical feature/observation versions, current
explanation members and their claim ancestors. Original observations required
by visible feature provenance are also retained. The query core starts from a
specified explanation, closes over declared claim parents, and retains cited
features resolved according to the query's temporal mode. Identity, container
and assessment objects become sidecars. Observation provenance is a separate
expansion. Claims' asserted intervals come from their proposition metadata,
not from the query-interval wrapper used by the runtime.

A citation remapped from an old immutable feature to its current logical version
is exported as CURRENT\_CITATION\_PATH with the original citation and version-path
witness IDs. It is never relabelled a directly stored Neo4j relationship.
Degree/depth calculations exclude SUPPORTS mirrors and count each dependency
edge once. Full storage counts immutable objects; core statistics count visible
propositions/versions per snapshot. Snapshots and replay copies remain clustered
within their source episode and participant.

The representative graph rule takes the lexicographically first median-size
eligible correction core per real source. Eligibility requires visible claims
and a revised cited feature. A separately registered maximum actual-impact rule
selects the same case in both sources, so no second purportedly distinct example
is drawn. The exports retain both rule labels and the complete eligible lists.
WESAD S11 and PPG-DaLiA S14 are therefore reproducible examples, not selected
because they show a propagation advantage. Both show direct operand citations.

\section{Independent Semantics and Extended Results}
\suppTable{generated/semantic/graph_table}{Node counts for full visible storage,
active evidence-and-claim graphs and query cores, with subject-cluster 95\%
intervals. Core/stored is a percentage; final column counts nonzero exposure
among eligible revision batches. Ten subjects and 20 episodes per source.}{tab:graphsizes}

The evaluator reads the observation/revision log and task specification; it never
imports graph support assessments or auxiliary-verifier decisions as reference
truth. Canonical identities include source, participant/session, interval,
quantity, operation, normalized unit/value, temporal mode and resolved logical
versions. Numerical matching preserves the original tolerance. Required slots
are matched once; raw duplicate propositions remain in the precision denominator.
Fully unanswerable requirements contribute NA recall, while partially answerable
tasks contribute only answerable slots. An empty answer to an answerable question
has zero recall. Subject means are computed after episode aggregation; pooled
counts are supplied for interpretation and are not the same estimand.

\suppTable{generated/semantic/crosswalk_table}{New semantic protocol versus the
unchanged B0 contract outcomes. Grounded background is a subset of grounded
formerly invalid content. The remaining disagreements include scope/dependency
and citation admission restrictions; no old valid claim becomes ungrounded.}{tab:crosswalk}
\suppTable{generated/semantic/competency_table}{Independent competency checks on
displayed M1 waveform claims. Each eligible comparison must expose both operands;
known-version and source-window answers are checked against immutable events.
Alternative-support behavior is tested separately in the structural experiment.}{tab:competency}
\suppTable{generated/semantic/temporal_table}{Event-aligned M1 prior-answer
checkpoints across all variants. Each source/position has 160 transitions from
20 episodes and ten subjects. Intervals are subject-cluster 95\% intervals.
No continuous-time loss or finite recovery time is inferred from these points.}{tab:temporal}

The machine-readable proposition ledger preserves each raw sentence, canonical
proposition, admissible numerical witnesses, match to requirements and limited
prose audit. That audit flags literal discrepancies and interpretation terms,
not every possible semantic implication. The new structural audit enumerates all
assignments to four or eight primitive tests and compares each proposed witness
with the task-defined proposition. Distinct logically equivalent derivations
are allowed. Soundness of each witness and completeness of the alternative
family are separate: a sufficient route can be sound while omitting another
valid route. True initial Boolean values alone do not establish either property.

\suppTable{generated/semantic/program_fidelity}{Structural candidate semantics:
initial propositions are true in all cases; witness soundness and answer-root coverage with
an admissible witness. Entries are percentages with subject-cluster intervals;
each source/stage has 80 cases from ten reused episodes and five subjects.
Accepted/displayed refers to original extension admission. Primitive test names
are normalized for independent proposition truth, not silently admitted.}{tab:programfidelity}

\section{Structural Registration, Admission Correction and Cost}

The initial protocol, deterministic selection and source hashes precede new
inference. Development pilot v1 produced terse test-token sentences and direct
roots. A second prospectively documented pilot clarified natural-language
sentences and reuse of earlier claim IDs, with a generic unrelated grouping
example. No required model links were compiled in. Both pilots use development
episodes; the core freeze follows pilot v2. A separate pilot-amendment manifest
records the added development budget and prompt hash. The core allows at most
one repair per case and no post-correction generation.

\suppTable{generated/semantic/extension_requests}{Separate extension inference
accounting. Roots count candidates admitted by that stage's runtime. All calls
use Qwen3-8B BF16 on the existing RTX 5090. CPU database replays add no model
requests. The core's 325 calls remain below its 480-call cap.}{tab:extensionrequests}

The frozen extension's proposition-name table omitted singleton primitive names
although their comparators were explicit in the input. The independent scorer
resolves such names to those definitions. A separately scoped corrected
replication additionally makes them admissible to the unchanged validator by
extending the proposition map. It preserves the final candidate JSON, source
events, witnesses, parent IDs, update conditions and repair opportunities.
Candidate hashes connect the corrected run to the original. It makes no new GPU
calls. Its \SemAliasRoots{} admitted roots replace neither the original
\SemCoreRoots{} admission result nor any minimum-v1 result.

\suppTable{generated/semantic/alias_table}{Separate primitive-name normalization
replication: 240 cases, four update conditions and five methods (4,800 cells).
All original model outputs and parent links are unchanged. Correction intervals
use 15 dataset-scoped participants and 30 episodes.}{tab:alias}

The normalized replication still realizes only evidence-to-answer depths one and
two. Undefined or malformed references are retained as failures; they are never
fixed by inserting edges. Its paired database outcomes agree, its eligible
direct-miss sets match exposure exactly, and its admitted two-route answers retain
support after primary-route loss. Full candidate/decision/replay records and
the independent audit are distributed alongside the primary run.

The structural replay timing fields include both direct and transitive reach
queries for diagnostic comparison, current snapshot retrieval, and evaluation
of the whole small program before scheduled writes. They are useful execution
accounting but not optimized affected-subgraph latency. Source-level numerical
corrections, benign changes and irrelevant changes are represented explicitly.
No price-per-hour claim is inferred because a billing tariff was not recorded.

\section{Preserved Minimum-v1 Accounting}

This section retains the completed study's definitions, counts and strict contract
interpretations. These are not retrospectively replaced by the new semantic
scores. The archived original main paper and report provide their full narrative.
Independent generation in the original five methods means conditional correction
denominators can differ; the new extension removes that ambiguity by replaying
one common candidate. The auxiliary verifier is a model audit, not a gold oracle.

\suppTable{generated/inventory_table}{Original source and prepared-record inventory.}{tab:inventory}
\suppTable{generated/method_table}{Original method ablations and maintenance behavior.}{tab:methods}
\suppTable{generated/reliability_table}{Frozen exact-contract reliability and
answer coverage. Correctly labelled off-target background remains a contract
error in this table.}{tab:reliability}
\suppTable{generated/correction_table}{Original required-change correction and
collateral effects, conditional on each method's displayed claims.}{tab:correction}
\suppTable{generated/contrasts_table}{Original paired descriptive method contrasts
with the original uncertainty procedure.}{tab:contrasts}
\suppTable{generated/fixture_table}{Original three-claim symbolic fixture,
kept separate from the new 240-case structural programs.}{tab:oldfixture}
\suppTable{generated/hardware_table}{Original hardware/model settings and sampled
resource peaks. Sampling can miss brief peaks.}{tab:hardware}
\suppTable{generated/failure_table}{Original failures, bounded repairs and
supported references. Three B0 transport errors remain recorded failures.}{tab:failures}
\suppTable{generated/contract_scope_table}{Original narrow background-scope audit.
These facts remain invalid under the frozen target-only contract.}{tab:oldscope}
\suppTable{generated/relational_contrasts_table}{Original graph versus relational
generated-outcome contrasts. Logical parity under common candidates is measured
separately in the new extension.}{tab:relational}
\suppTable{generated/streaming_table}{Original systems events, backlog and latency
accounting. Finite candidate-completion latency does not represent adequate
replacement time.}{tab:streaming}
\suppTable{generated/semantic/cost_table}{Flagging and candidate completion versus
adequate replacement: one 60-second cell per offered rate and method.
All successful-replacement times remain unresolved.}{tab:compactcost}
\begin{figure}[!htbp]
\centering\includegraphics[width=\textwidth]{streaming.pdf}
\caption{Preserved four-panel systems measurements. Finite generation latency
is candidate-completion latency. No candidate satisfies both required facts,
so these curves do not measure successful correction time.}
\end{figure}
\suppTable{generated/staleness_table}{Completion-time adequacy audit. Zero adequate
answers makes the zero became-stale count uninformative about robust recovery.}{tab:staleness}
\suppTable{generated/cost_table}{Original request/token accounting by workload.
Extension costs appear separately in Table~\ref{tab:extensionrequests}.}{tab:oldcost}

\clearpage
\section{Reproduction and Artifact Map}

From the project root, the existing virtual environment and downloaded prepared
events suffice for CPU analysis; no minimum-v1 model calls are repeated.
The commands and their complete options are maintained in
\path{docs/semantic_reproduction.md}. Run the semantic modules in order:
analysis, structural analysis, alias analysis, figures and publication.
Their exact module names and commands are listed in that document.
Then run \texttt{bash paper/build.sh}.
Focused tests and the PDF/integrity checker are included.

Actual original database exports reside in
\path{artifacts/analysis/semantic_analysis_v1/exports/};
the matching manifest preserves Cypher, scope, export time, checksums and
representative stored paths. Selected structural database exports live under
\texttt{structure\_exports/}. Semantic projection sidecars and figure layout
files retain full IDs and derived-path witnesses. The original inputs and runtime
hashes are verified against the minimum-run manifest. New inference uses
\texttt{artifacts/runs/structure\_study\_v1/}, with core prompts, raw responses,
repairs, CUDA placement checks, request journals, resource samples and all
method/update logs. The separate normalization replication has its own
manifest, candidates and replay directory.

This Supplement and the \href{semantic_structure_revision.pdf}{main manuscript}
use the official LLNCS template. The final report records
the actual page counts, figure checks and source/result hashes. No conference
submission is performed by the reproduction workflow.
\end{document}
